\documentclass[11pt]{article}
\usepackage{amsmath,amssymb,geometry}
\geometry{margin=1in}
\title{ROOT V151: Helicity and Momentum-Space Berry Geometry (Audit)}
\date{24 November 2026}
\begin{document}\maketitle
\noindent\textbf{Unit convention:} only $c=1$; retain $\hbar$ and $G$ with dimensions.
\section*{Moving orthogonal frame}
\begin{align}
\mathbf n&=(\sin\theta\cos\phi,\sin\theta\sin\phi,\cos\theta),\\
\mathbf e_\theta&=(\cos\theta\cos\phi,\cos\theta\sin\phi,-\sin\theta),\\
\mathbf e_\phi&=(-\sin\phi,\cos\phi,0),\\
\mathbf e_\sigma&=(\mathbf e_\theta+i\sigma\mathbf e_\phi)/\sqrt2,\quad \sigma=\pm1.
\end{align}
\section*{Photon helicity and Berry curvature}
\begin{align}
(i\mathbf n\times)\mathbf e_\sigma&=\sigma\mathbf e_\sigma,\\
\mathcal A_\sigma&=-i\mathbf e_\sigma^\dagger d\mathbf e_\sigma=-\sigma\cos\theta\,d\phi,\\
\mathcal F_\sigma&=d\mathcal A_\sigma=\sigma\sin\theta\,d\theta\wedge d\phi,\\
C_\sigma&=\frac{1}{2\pi}\int_{S^2}\mathcal F_\sigma=2\sigma,\\
\boldsymbol\Omega_\sigma(\mathbf k)&=\sigma\mathbf k/|\mathbf k|^3,\qquad\mathbf k\ne0.
\end{align}
\noindent This Berry curvature lives in momentum-direction space, not in spacetime.
\section*{Gauge redundancy and Maxwell}
\begin{align}
\mathbf e_\sigma &\to e^{i\chi}\mathbf e_\sigma, &\mathcal A&\to\mathcal A+d\chi,
&\mathcal F&\to\mathcal F,\\
A_\mu&=\partial_\mu\psi &\Rightarrow F_{\mu\nu}&=\partial_\mu A_\nu-\partial_\nu A_\mu=0.
\end{align}
\section*{Circular polarization}
\begin{align}
\varphi&=\omega(t-z),\\
\mathbf E&=E_0(\cos\varphi,\sigma\sin\varphi,0),\\
\mathbf B&=E_0(-\sigma\sin\varphi,\cos\varphi,0),\\
\nabla\cdot\mathbf E&=\nabla\cdot\mathbf B=0,\\
\partial_t\mathbf E&=\nabla\times\mathbf B,\qquad \partial_t\mathbf B=-\nabla\times\mathbf E,\\
\mathbf E\cdot\mathbf B&=0,\quad E^2=B^2,\quad\mathbf E\times\mathbf B=E_0^2\hat{\mathbf z}.
\end{align}
\section*{Lorentz transformation}
\begin{equation}
D=\gamma(1-\beta)=\sqrt{\frac{1-\beta}{1+\beta}},\qquad
\omega'=D\omega,\quad E_0'=DE_0,\quad E_\gamma'=\hbar\omega'.
\end{equation}
\section*{Limits of the claim}
Neither $C_\sigma=\pm2$ nor an orthogonal three-frame derives
$\alpha=e^2/(4\pi\epsilon_0\hbar)$ (for $c=1$), $G$, or the Standard Model
internal gauge group.  The equations here are established mathematical
identities and Maxwell-model tests, not a proof of four-force unification.
\end{document}
